%!TEX root=./main.tex

\section{Related Work}

In this paper, we have studied the problem of estimating the number of Mazurkiewicz traces in a program. 
% \paragraph{Approximate Counting}{
(Approximate) Counting of combinatorial objects is a classical topic in complexity theory \cite{Valiant79,JerrumVV86}.
A central problem here is \emph{propositional model counting} ($\#SAT$).
Our proof of hardness of the estimation problem reduces from $\#SAT$.
There are approximate counters for $\#SAT$, such as ApproxMC \cite{chakraborty2013scalableapproximatemodelcounter}, that run in randomized polynomial time with access to
an NP-oracle.
These techniques are not easily applied in our context because compiling all executions of a real-world concurrent program into a SAT instance is a formidable problem.
Instead, we take the route of unbiased estimators obtained from stateless search.

Our work is also close to approximate counting on automata, where one asks for the number of words of a given length accepted by an NFA. 
This problem admits a fully polynomial randomized approximation scheme (FPRAS) \cite{GoreJKSM97,10.1145/3422648.3422661}, 
with several algorithmic improvements \cite{meel2024fasterfprasnfa,meel2025practicalfprasnfaexploiting}.
A recent paper \cite{ColnetMM26} considers the problem of counting traces in a regular trace language and shows an FPRAS for the problem.
It might seem that our estimation problem can be reduced to these problems by taking the automaton for each thread, forming their partially synchronous product, 
intersecting with the regular language of lexicographic normal forms of trace equivalence, and finally asking how many words of some appropriate length the resulting automaton accepts. 
However, the product automaton is exponential in the number of threads and variables, so even FPRAS algorithms for approximate counting on automata become exponential in the number of threads. 
%% Moreover, even finding the automata for each thread of a concurrent program statically is a difficult problem.

%Similar problems have been studied in the literature, albeit with different approaches. 
Our work is also related to estimating coverage in a model checker.
Penix and Visser~\cite{penix} add standard code-coverage metrics to Java PathFinder \cite{jpf}, tracking which bytecode instructions and branch outcomes are exercised during state-space exploration.
Bron, Ziv, and Ur~\cite{bron05} introduce synchronization coverage for concurrent Java programs, defining a small set of concurrency-specific coverage tasks (for example, whether a synchronized block ever caused another thread to block) and treating full coverage as a practical stopping criterion.
Nevertheless, code coverage and its derivatives are known to be imprecise measures for the space of a program.
% and note that the resulting coverage curves often “converge’’ once further exploration rarely reveals new program elements. 
Chockler, Kupferman, and Vardi~\cite{10.1007/978-3-540-39724-3_11} propose mutation-based coverage metrics for model checking that ask which states, transitions, and subformulas actually matter for establishing
if a specification is satisfied.  
% An element is considered covered if mutating it falsifies the property or makes it hold vacuously, and they give symbolic algorithms instantiating this idea for several standard coverage notions. 
Taleghani and Atlee \cite{5431751} estimate state-space coverage in explicit-state model checking by sampling unexplored transitions using BFS and extrapolating the number of remaining states. 
However the estimate is heuristic and comes with no guarantees of being an unbiased estimator. 
% Path profiling proposed by Ball and Larus \cite{10.5555/243846.243857} is another dynamic coverage metric: instead of asking which blocks or edges are covered, it measures how often each control-flow path in a procedure is executed under a given workload. 
% The classic scheme of Ball and Larus also relies on a DAG view of executions, obtained by cutting loop back-edges in the control-flow graph, and uses this structure to profile path frequencies, which shares certain similar sprit of our DAG estimator.
These approaches describe which parts of the system and specification are exercised, or roughly how much of the explicit state graph has been visited, but not how far verification has progressed within the underlying execution space. 
In contrast, our contribution is complementary: we start with a complete set of behaviors as our coverage goal and we estimate the size of the set.
Moreover, our notion of behavior identifies two causally identical executions as one.
Closest to our work is the estimate of \cite{gator}, which inspired this work.
As discussed before, our paper provides the precise complexity bounds for the estimation problem and ensures that the estimator is unbiased.

Purely statistical techniques have also been proposed in recent  fuzzing literature.
For example, the use of Good-Turing estimator \cite{10.1145/3210309, 10.1145/3611019,10.1145/3468264.3468570} has been proposed to give statistical evidence to continue or to stop a fuzzing campaign; it is unclear if their estimate can
be converted to an estimate for the concurrent state space.
More heuristic notions of saturation or past coverage trends have also been proposed to predict the expected time to full coverage  \cite{10.1145/3597503.3639198,10.1145/3597926.3598043}. 
% As such, these problems are specific to purely memoryless testing techniques like greybox and blackbox fuzzing, that aim to discover as many bugs as possible.
% Several works in fuzzing propose estimators to reason about how much testing has achieved. The STADS framework and Entropic power schedules model fuzzing as a species-discovery process and adapt Good–Turing–style estimators to predict, for example, the probability that the next test input will exercise a new behaviour and the eventual number of behaviours a fuzzer might discover~\cite{10.1145/3210309, 10.1145/3611019}. Subsequent work extrapolates future coverage from past coverage curves or uses the saturation of predicted-vulnerable functions as a heuristic stopping rule. These approaches are primarily evaluated empirically and, in the greybox setting, do not come with provable unbiasedness or variance guarantees. 
%
% Our setting is different: we aim to count semantic execution equivalence classes of a concurrent program, and we design estimators with explicit mathematical guarantees for this counting problem, rather than heuristic progress indicators for fuzzing campaigns.





\section{Conclusions and Future Work}

In this paper, we looked at the problem of estimating the size of the space explored by a model checker on a bounded concurrent program. Even though there has been previous work on related questions, they were based on heuristic approaches without any theoretical results. 
We lay the theoretical foundations for this problem, both by exhibiting hardness results as well as by coming up with the first poly-time unbiased estimator for the problem. 
An implementation of our procedure on a wide range of shared-memory concurrency benchmarks shows that our estimator provides robust estimates whilst only using a modest number of samples.
This work  opens up avenues for interesting future work including, accounting 
for weaker equivalences such as reads-from equivalence~\cite{readsFrom2019},
reads-value-from equivalences~\cite{readsValue2021} and recent intermediate
equivalences such as grain equivalence~\cite{FarzanMathur2024} as well as weaker memory models.

 
